% Diagrama de flujo CQRS / Event Sourcing
% Uso: % Diagrama de flujo CQRS / Event Sourcing
% Uso: % Diagrama de flujo CQRS / Event Sourcing
% Uso: % Diagrama de flujo CQRS / Event Sourcing
% Uso: \input{figs/cqrs-flujo} dentro de figure+\centering+\resizebox{\textwidth}{!}{...}
%
% Flujo: Comando → Manejador → Agregado → Evento → EventStore → Proyector → Modelo de lectura

\begin{tikzpicture}[
  every node/.style={font=\small\sffamily},
  %--- Escritura (lado izquierdo) ---
  cmd/.style={
    draw=palered!80, fill=palered!10, thick,
    rectangle, rounded corners=4pt,
    minimum width=2.6cm, minimum height=1.1cm,
    align=center
  },
  handler/.style={
    draw=gray!60, fill=gray!10, thick,
    rectangle, rounded corners=4pt,
    minimum width=2.6cm, minimum height=1.1cm,
    align=center
  },
  agg/.style={
    draw=aquaESI!80, fill=aquaESI!15, thick,
    rectangle, rounded corners=4pt,
    minimum width=2.6cm, minimum height=1.1cm,
    align=center
  },
  evt/.style={
    draw=ocre!80, fill=ocre!15, thick,
    rectangle, rounded corners=4pt,
    minimum width=2.6cm, minimum height=1.1cm,
    align=center
  },
  store/.style={
    draw=turquesa!80, fill=turquesa!15, thick,
    cylinder, shape border rotate=90,
    minimum height=1.4cm, minimum width=2.4cm,
    align=center
  },
  proj/.style={
    draw=gray!70, fill=gray!12, thick,
    rectangle, rounded corners=4pt,
    minimum width=2.6cm, minimum height=1.1cm,
    align=center
  },
  read/.style={
    draw=aquaESI!60, fill=aquaESI!10, thick,
    rectangle, rounded corners=4pt,
    minimum width=2.6cm, minimum height=1.1cm,
    align=center
  },
  %--- Flechas ---
  arr/.style={->, >=Stealth, thick},
  lbl/.style={font=\small\sffamily, midway, fill=white,
              inner sep=1.5pt, draw=gray!25, rounded corners=2pt},
  node distance=0.5cm and 0.9cm
]

%--- LADO ESCRITURA (fila superior) ---
\node[cmd] (cmd1) at (0,0) {
  \textbf{Comando}\\
  \footnotesize AnalyzeReport,\\ValidateCode
};

\node[agg, right=2.2cm of cmd1] (agg) {
  \textbf{Agregado}\\
  \footnotesize Analysis\\(estado)
};

\node[evt, right=2.0cm of agg] (evt1) {
  \textbf{Evento}\\
  \footnotesize ReportAnalyzed,\\CodeValidated
};

\node[store, right=2.2cm of evt1] (es) {
  \textbf{Event Store}\\
  \footnotesize PostgreSQL
};

%--- PROYECTOR (fila inferior) ---
\node[proj, below=1.4cm of es] (proj) {
  \textbf{Proyector}\\
  \footnotesize una proyección\\por evento
};

\node[read, left=1.2cm of proj, minimum width=5.4cm, text width=5.2cm] (r1) {
  \textbf{Modelo de lectura}\\
  \footnotesize conversations, analysis\_cards,\\predicted\_codes
};

%--- FLECHAS DE ESCRITURA ---
\draw[arr] (cmd1) -- node[lbl, above]{ejecuta} (agg);
\draw[arr] (agg)  -- node[lbl, above]{emite} (evt1);
\draw[arr] (evt1) -- node[lbl, above]{guarda} (es);

%--- FLECHA EventStore → Proyector ---
\draw[arr] (es) -- node[lbl, right]{se suscribe} (proj);

\draw[arr] (proj) -- node[lbl, above]{actualiza} (r1);

%--- ETIQUETAS DE ZONA ---
\node[font=\small\sffamily\bfseries, palered!80]
  at ($(cmd1)!0.5!(evt1) + (0,1.5)$) {Lado escritura (comandos)};

\node[font=\small\sffamily\bfseries, aquaESI!80]
  at ($(r1)!0.5!(proj) + (0,-1.5)$) {Lado lectura (consultas y proyecciones)};

%--- SEPARADOR ---
\draw[dashed, gray!50, thick]
  ($(es.south)!0.5!(proj.north)$) ++(-10,0) -- ++(13,0);

\end{tikzpicture}
 dentro de figure+\centering+\resizebox{\textwidth}{!}{...}
%
% Flujo: Comando → Manejador → Agregado → Evento → EventStore → Proyector → Modelo de lectura

\begin{tikzpicture}[
  every node/.style={font=\small\sffamily},
  %--- Escritura (lado izquierdo) ---
  cmd/.style={
    draw=palered!80, fill=palered!10, thick,
    rectangle, rounded corners=4pt,
    minimum width=2.6cm, minimum height=1.1cm,
    align=center
  },
  handler/.style={
    draw=gray!60, fill=gray!10, thick,
    rectangle, rounded corners=4pt,
    minimum width=2.6cm, minimum height=1.1cm,
    align=center
  },
  agg/.style={
    draw=aquaESI!80, fill=aquaESI!15, thick,
    rectangle, rounded corners=4pt,
    minimum width=2.6cm, minimum height=1.1cm,
    align=center
  },
  evt/.style={
    draw=ocre!80, fill=ocre!15, thick,
    rectangle, rounded corners=4pt,
    minimum width=2.6cm, minimum height=1.1cm,
    align=center
  },
  store/.style={
    draw=turquesa!80, fill=turquesa!15, thick,
    cylinder, shape border rotate=90,
    minimum height=1.4cm, minimum width=2.4cm,
    align=center
  },
  proj/.style={
    draw=gray!70, fill=gray!12, thick,
    rectangle, rounded corners=4pt,
    minimum width=2.6cm, minimum height=1.1cm,
    align=center
  },
  read/.style={
    draw=aquaESI!60, fill=aquaESI!10, thick,
    rectangle, rounded corners=4pt,
    minimum width=2.6cm, minimum height=1.1cm,
    align=center
  },
  %--- Flechas ---
  arr/.style={->, >=Stealth, thick},
  lbl/.style={font=\small\sffamily, midway, fill=white,
              inner sep=1.5pt, draw=gray!25, rounded corners=2pt},
  node distance=0.5cm and 0.9cm
]

%--- LADO ESCRITURA (fila superior) ---
\node[cmd] (cmd1) at (0,0) {
  \textbf{Comando}\\
  \footnotesize AnalyzeReport,\\ValidateCode
};

\node[agg, right=2.2cm of cmd1] (agg) {
  \textbf{Agregado}\\
  \footnotesize Analysis\\(estado)
};

\node[evt, right=2.0cm of agg] (evt1) {
  \textbf{Evento}\\
  \footnotesize ReportAnalyzed,\\CodeValidated
};

\node[store, right=2.2cm of evt1] (es) {
  \textbf{Event Store}\\
  \footnotesize PostgreSQL
};

%--- PROYECTOR (fila inferior) ---
\node[proj, below=1.4cm of es] (proj) {
  \textbf{Proyector}\\
  \footnotesize una proyección\\por evento
};

\node[read, left=1.2cm of proj, minimum width=5.4cm, text width=5.2cm] (r1) {
  \textbf{Modelo de lectura}\\
  \footnotesize conversations, analysis\_cards,\\predicted\_codes
};

%--- FLECHAS DE ESCRITURA ---
\draw[arr] (cmd1) -- node[lbl, above]{ejecuta} (agg);
\draw[arr] (agg)  -- node[lbl, above]{emite} (evt1);
\draw[arr] (evt1) -- node[lbl, above]{guarda} (es);

%--- FLECHA EventStore → Proyector ---
\draw[arr] (es) -- node[lbl, right]{se suscribe} (proj);

\draw[arr] (proj) -- node[lbl, above]{actualiza} (r1);

%--- ETIQUETAS DE ZONA ---
\node[font=\small\sffamily\bfseries, palered!80]
  at ($(cmd1)!0.5!(evt1) + (0,1.5)$) {Lado escritura (comandos)};

\node[font=\small\sffamily\bfseries, aquaESI!80]
  at ($(r1)!0.5!(proj) + (0,-1.5)$) {Lado lectura (consultas y proyecciones)};

%--- SEPARADOR ---
\draw[dashed, gray!50, thick]
  ($(es.south)!0.5!(proj.north)$) ++(-10,0) -- ++(13,0);

\end{tikzpicture}
 dentro de figure+\centering+\resizebox{\textwidth}{!}{...}
%
% Flujo: Comando → Manejador → Agregado → Evento → EventStore → Proyector → Modelo de lectura

\begin{tikzpicture}[
  every node/.style={font=\small\sffamily},
  %--- Escritura (lado izquierdo) ---
  cmd/.style={
    draw=palered!80, fill=palered!10, thick,
    rectangle, rounded corners=4pt,
    minimum width=2.6cm, minimum height=1.1cm,
    align=center
  },
  handler/.style={
    draw=gray!60, fill=gray!10, thick,
    rectangle, rounded corners=4pt,
    minimum width=2.6cm, minimum height=1.1cm,
    align=center
  },
  agg/.style={
    draw=aquaESI!80, fill=aquaESI!15, thick,
    rectangle, rounded corners=4pt,
    minimum width=2.6cm, minimum height=1.1cm,
    align=center
  },
  evt/.style={
    draw=ocre!80, fill=ocre!15, thick,
    rectangle, rounded corners=4pt,
    minimum width=2.6cm, minimum height=1.1cm,
    align=center
  },
  store/.style={
    draw=turquesa!80, fill=turquesa!15, thick,
    cylinder, shape border rotate=90,
    minimum height=1.4cm, minimum width=2.4cm,
    align=center
  },
  proj/.style={
    draw=gray!70, fill=gray!12, thick,
    rectangle, rounded corners=4pt,
    minimum width=2.6cm, minimum height=1.1cm,
    align=center
  },
  read/.style={
    draw=aquaESI!60, fill=aquaESI!10, thick,
    rectangle, rounded corners=4pt,
    minimum width=2.6cm, minimum height=1.1cm,
    align=center
  },
  %--- Flechas ---
  arr/.style={->, >=Stealth, thick},
  lbl/.style={font=\small\sffamily, midway, fill=white,
              inner sep=1.5pt, draw=gray!25, rounded corners=2pt},
  node distance=0.5cm and 0.9cm
]

%--- LADO ESCRITURA (fila superior) ---
\node[cmd] (cmd1) at (0,0) {
  \textbf{Comando}\\
  \footnotesize AnalyzeReport,\\ValidateCode
};

\node[agg, right=2.2cm of cmd1] (agg) {
  \textbf{Agregado}\\
  \footnotesize Analysis\\(estado)
};

\node[evt, right=2.0cm of agg] (evt1) {
  \textbf{Evento}\\
  \footnotesize ReportAnalyzed,\\CodeValidated
};

\node[store, right=2.2cm of evt1] (es) {
  \textbf{Event Store}\\
  \footnotesize PostgreSQL
};

%--- PROYECTOR (fila inferior) ---
\node[proj, below=1.4cm of es] (proj) {
  \textbf{Proyector}\\
  \footnotesize una proyección\\por evento
};

\node[read, left=1.2cm of proj, minimum width=5.4cm, text width=5.2cm] (r1) {
  \textbf{Modelo de lectura}\\
  \footnotesize conversations, analysis\_cards,\\predicted\_codes
};

%--- FLECHAS DE ESCRITURA ---
\draw[arr] (cmd1) -- node[lbl, above]{ejecuta} (agg);
\draw[arr] (agg)  -- node[lbl, above]{emite} (evt1);
\draw[arr] (evt1) -- node[lbl, above]{guarda} (es);

%--- FLECHA EventStore → Proyector ---
\draw[arr] (es) -- node[lbl, right]{se suscribe} (proj);

\draw[arr] (proj) -- node[lbl, above]{actualiza} (r1);

%--- ETIQUETAS DE ZONA ---
\node[font=\small\sffamily\bfseries, palered!80]
  at ($(cmd1)!0.5!(evt1) + (0,1.5)$) {Lado escritura (comandos)};

\node[font=\small\sffamily\bfseries, aquaESI!80]
  at ($(r1)!0.5!(proj) + (0,-1.5)$) {Lado lectura (consultas y proyecciones)};

%--- SEPARADOR ---
\draw[dashed, gray!50, thick]
  ($(es.south)!0.5!(proj.north)$) ++(-10,0) -- ++(13,0);

\end{tikzpicture}
 dentro de figure+\centering+\resizebox{\textwidth}{!}{...}
%
% Flujo: Comando → Manejador → Agregado → Evento → EventStore → Proyector → Modelo de lectura

\begin{tikzpicture}[
  every node/.style={font=\small\sffamily},
  %--- Escritura (lado izquierdo) ---
  cmd/.style={
    draw=palered!80, fill=palered!10, thick,
    rectangle, rounded corners=4pt,
    minimum width=2.6cm, minimum height=1.1cm,
    align=center
  },
  handler/.style={
    draw=gray!60, fill=gray!10, thick,
    rectangle, rounded corners=4pt,
    minimum width=2.6cm, minimum height=1.1cm,
    align=center
  },
  agg/.style={
    draw=aquaESI!80, fill=aquaESI!15, thick,
    rectangle, rounded corners=4pt,
    minimum width=2.6cm, minimum height=1.1cm,
    align=center
  },
  evt/.style={
    draw=ocre!80, fill=ocre!15, thick,
    rectangle, rounded corners=4pt,
    minimum width=2.6cm, minimum height=1.1cm,
    align=center
  },
  store/.style={
    draw=turquesa!80, fill=turquesa!15, thick,
    cylinder, shape border rotate=90,
    minimum height=1.4cm, minimum width=2.4cm,
    align=center
  },
  proj/.style={
    draw=gray!70, fill=gray!12, thick,
    rectangle, rounded corners=4pt,
    minimum width=2.6cm, minimum height=1.1cm,
    align=center
  },
  read/.style={
    draw=aquaESI!60, fill=aquaESI!10, thick,
    rectangle, rounded corners=4pt,
    minimum width=2.6cm, minimum height=1.1cm,
    align=center
  },
  %--- Flechas ---
  arr/.style={->, >=Stealth, thick},
  lbl/.style={font=\small\sffamily, midway, fill=white,
              inner sep=1.5pt, draw=gray!25, rounded corners=2pt},
  node distance=0.5cm and 0.9cm
]

%--- LADO ESCRITURA (fila superior) ---
\node[cmd] (cmd1) at (0,0) {
  \textbf{Comando}\\
  \footnotesize AnalyzeReport,\\ValidateCode
};

\node[agg, right=2.2cm of cmd1] (agg) {
  \textbf{Agregado}\\
  \footnotesize Analysis\\(estado)
};

\node[evt, right=2.0cm of agg] (evt1) {
  \textbf{Evento}\\
  \footnotesize ReportAnalyzed,\\CodeValidated
};

\node[store, right=2.2cm of evt1] (es) {
  \textbf{Event Store}\\
  \footnotesize PostgreSQL
};

%--- PROYECTOR (fila inferior) ---
\node[proj, below=1.4cm of es] (proj) {
  \textbf{Proyector}\\
  \footnotesize una proyección\\por evento
};

\node[read, left=1.2cm of proj, minimum width=5.4cm, text width=5.2cm] (r1) {
  \textbf{Modelo de lectura}\\
  \footnotesize conversations, analysis\_cards,\\predicted\_codes
};

%--- FLECHAS DE ESCRITURA ---
\draw[arr] (cmd1) -- node[lbl, above]{ejecuta} (agg);
\draw[arr] (agg)  -- node[lbl, above]{emite} (evt1);
\draw[arr] (evt1) -- node[lbl, above]{guarda} (es);

%--- FLECHA EventStore → Proyector ---
\draw[arr] (es) -- node[lbl, right]{se suscribe} (proj);

\draw[arr] (proj) -- node[lbl, above]{actualiza} (r1);

%--- ETIQUETAS DE ZONA ---
\node[font=\small\sffamily\bfseries, palered!80]
  at ($(cmd1)!0.5!(evt1) + (0,1.5)$) {Lado escritura (comandos)};

\node[font=\small\sffamily\bfseries, aquaESI!80]
  at ($(r1)!0.5!(proj) + (0,-1.5)$) {Lado lectura (consultas y proyecciones)};

%--- SEPARADOR ---
\draw[dashed, gray!50, thick]
  ($(es.south)!0.5!(proj.north)$) ++(-10,0) -- ++(13,0);

\end{tikzpicture}
